% This file is intended to be \input{} from subsubsec__window6_lie_bracket_closure.tex.

\paragraph{交换子包络、根字典与残余结构：从离散审计到可验证闭合链}
本段将 window-\(6\) 处“李括号的构造”与“唯一闭合”的两条对象链作同一口径下的衔接：其一是由坐标翻转推送算子生成的交换子 Lie 包络（算子链），其二是由根--Cartan 字典在 \(V=\RR\langle X_6\rangle\) 上诱导的抽象括号（类型链）。
两条链均以有限可审计数据为输入，但其闭合目标不同；因此相关猜想的定理化应当分别落在这两条对象链的相应命题上。

\begin{remark}[交换子包络的严格构造与唯一闭合]\label{rem:window6-commutator-envelope-unique-closure}
对每个坐标方向 \(i\in\{1,\dots,6\}\)，推送计数算子 \(L_i\) 由有限边计数确定（定义 \ref{def:window6-coordinate-pushforward-lie-envelope}），并以交换子 \([A,B]=AB-BA\) 作为李括号。
由最小交刻画，交换子包络
\(
\mathfrak g_6=\mathrm{Lie}\langle L_1,\dots,L_6\rangle
\)
作为包含生成元集合的最小 Lie 子代数存在且唯一（定理 \ref{thm:window6-lie-envelope-existence-uniqueness}），并可通过有限维闭包迭代在有限步内稳定（命题 \ref{prop:finite-dim-lie-closure-termination}）。
这些结论仅依赖有限维线性代数与交换子封闭性，不引入任何连续假设。
\end{remark}

\begin{remark}[关于猜想 \(10.170\) 的对象更正]\label{rem:window6-conj-10170-object-correction}
在同一交换子口径下，\(\mathfrak g_6\) 的结构识别已由证书化审计闭合为满的 \(\mathfrak{sl}_{21}\)（定理 \ref{thm:window6-lie-envelope-sl21}），从而排除“\(\dim\mathfrak g_6=21\) 且根型为 \(B_3/C_3\)”这一对交换子包络本身的预期。
因此，秩 \(3\) 的 \(21\) 维简单根型结构应当被理解为类型链上的抽象括号对象（定理 \ref{thm:window6-root-dictionary-pullback-bracket} 与猜想 \ref{conj:window6-intrinsic-bracket-uniqueness}），或作为 \(\mathfrak{sl}_{21}\) 内受限可观测紧致部分的抽取目标（猜想 \ref{conj:window6-sm-gauge-as-minimal-compact-observable}），而非 \(\mathfrak g_6\) 的交换子闭包结构本身。
\end{remark}

\begin{remark}[两类分块骨架与 SM 残余结构的兼容闭合]\label{rem:window6-two-skeletons-to-sm-residual}
window-\(6\) 同时提供两条彼此独立的离散分块接口：其一是 edge--flux 口径下的四块分解 \(21=9\oplus 6\oplus 3\oplus 3\)（终端命题 \ref{thm:terminal-window6-edge-flux-skeleton}），其二是 orbit--Levi 口径下的 \(21=15\oplus 3\oplus 3\) 维数骨架及相应 Levi 候选
\(
\mathfrak g^{\mathrm{Levi}}_6\cong \mathfrak{su}(4)\oplus\mathfrak{su}(2)_L\oplus\mathfrak{su}(2)_R
\)
（参见 \S\ref{subsec:bdry-tower-zeck-gut} 与命题 \ref{prop:window6-common-refinement-sm-levi}）。
当要求两条分块在同一连续包络/可观测语义下兼容闭合时，标准模型型 Levi 子结构
\(
\mathfrak{su}(3)\oplus\mathfrak{su}(2)_L\oplus\mathfrak u(1)
\)
被强制出现（命题 \ref{prop:window6-common-refinement-sm-levi}）。
进一步地，在 Levi 骨架之上，\(\mathfrak{su}(4)\supset\mathfrak{su}(3)\oplus\mathfrak u(1)\) 的嵌入刚性（推论 \ref{cor:window6-su4-embedding-rigidity}）与几何 \(\ZZ_2\)/sheet parity 审计共同给出一条有效阿贝尔方向 \(\mathfrak u(1)_{\mathrm{eff}}\subset\mathfrak{su}(2)_R\)，从而得到标准模型型残余代数（推论 \ref{cor:window6-levi-rigidity-sm-residual}）。
\end{remark}

\begin{remark}[抽象括号的存在性、唯一性与混合证书化路径]\label{rem:window6-bracket-theoremization-route}
根--Cartan 字典在 \(V=\RR\langle X_6\rangle\) 上给出一条完全显式的括号存在性构造：将 \(X_6^{\mathrm{bdry}}\) 锚定到 Cartan 基、将 \(X_6^{\mathrm{cyc}}\) 通过根字典锚定到根向量，从而拉回紧致 \(B_3/C_3\) 型括号（定理 \ref{thm:window6-root-dictionary-pullback-bracket}）。
在固定标号刚性、分级相容与分块相容的审计语义下，其唯一性可被组织为证书化断言（定理 \ref{thm:window6-intrinsic-bracket-certificate-theorem}）。
\par
在该框架中，括号的“存在/唯一”可进一步增强为混合证书问题：代数侧以反对称性与 Jacobi 恒等式、局域支撑模板及 parity/分块约束给出有限多项式系统（命题 \ref{prop:window6-bracket-certificate-reduction}），谱侧以 unitary slice 的 Toeplitz--PSD 证书链给出半正定约束（定理 \ref{thm:window6-p6-toeplitz-certificate-chain}、定理 \ref{thm:app-horizon-toeplitz-lmi}，并见定理 \ref{thm:window6-bracket-uniqueness-mixed-certificate-reduction}）。
由此，原先的“纯 Jacobi 多项式系统”可被替换为“多项式可行域 \(\cap\) 半正定可行域”的交汇判别，从而对非紧/非幺正分支给出可复核的排除机制。
\end{remark}

\begin{remark}[将猜想 \(10.170\)–\(10.172\) 升格为定理的缺口定位]\label{rem:window6-10-170-10-172-gaps}
交换子包络链的闭合与结构识别已在定理 \ref{thm:window6-lie-envelope-sl21} 中完成证书化。
剩余的定理化工作集中在两处：其一是把“可观测极小紧致部分”的存在与极小性准则给出为可审计判别（对应猜想 \ref{conj:window6-sm-gauge-as-minimal-compact-observable} 的实质内容）；其二是在类型链上把抽象括号的审计约束（局域/分级/分块）与混合证书约束（Toeplitz--PSD 等）合并为可输出证书的唯一可解性结论（对应猜想 \ref{conj:window6-intrinsic-bracket-uniqueness} 的定理化版本）。
两处缺口均可被组织为有限对象的核验任务，并可与整数模素数审计与 Lee--Yang/PSD 证书链并行交叉复核。
\end{remark}

